% 图 55.1 —— 由 `checks/emit_hand.py` 自动拼出（probe 精测 + prims2 图元分解）
%%
% ⚠ 这是**稿子不是成品**：跑 `checks/checkhand.py 55.1` 看指标 + 看叠合图。
%%
%% ⚠ 待核：
%   · 本图 probe 报出阴影族 1 个 —— **本脚本不生成阴影**（要照原书手写）；prims2 的 strokes 里可能含阴影线，核对时留意别画重
%   · 可疑字形 1 项（空心圈 ⊙ 常被认成字母 o）—— 见文件末
%%
\documentclass[12pt]{standalone}
\usepackage{fontspec}
\setsansfont{Arial}
\newfontfamily\mfallback{DejaVu Sans}
\usepackage{tikz}
\begin{document}
\begin{tikzpicture}[x=1pt, y=-1pt, line width=4.0pt, line cap=round, line join=round]

  %% ════ 填充区（prims2 的 fills）════

  %% ════ 直线段（prims2 的 strokes）════
  \draw[line width=5.0pt] (763.0,195.0) -- (787.0,169.0);
  \draw[line width=5.0pt] (762.0,196.0) -- (766.4,214.4);
  \draw[line width=5.0pt] (766.4,214.4) -- (775.8,220.8);
  \draw[line width=5.0pt] (897.8,184.2) -- (898.0,174.1);
  \draw[line width=5.0pt] (106.0,130.0) -- (58.0,120.0);
  \draw[line width=7.0pt] (710.0,88.0) -- (694.0,70.0);
  \draw[line width=4.0pt] (336.0,87.0) -- (471.0,87.0);
  \draw[line width=6.1pt] (469.0,81.0) -- (472.0,86.0);
  \draw[line width=4.0pt] (38.0,110.0) -- (36.0,171.0);
  \draw[line width=7.0pt] (473.0,87.0) -- (490.0,86.0);
  \draw[line width=5.0pt] (108.0,130.0) -- (162.0,124.0);
  \draw[line width=5.0pt] (818.0,115.0) -- (862.0,99.0);
  \draw[line width=4.0pt] (435.0,314.0) -- (508.0,365.0);
  \draw[line width=5.9pt] (710.0,89.0) -- (711.4,94.4);
  \draw[line width=5.8pt] (758.0,190.8) -- (761.0,195.0);
  \draw[line width=5.0pt] (215.0,95.2) -- (213.0,90.0);
  \draw[line width=3.0pt] (708.0,70.0) -- (710.2,68.0);
  \draw[line width=5.8pt] (859.2,94.2) -- (863.0,98.0);
  \draw[line width=9.0pt] (708.0,70.0) -- (695.0,69.0);
  \draw[line width=4.9pt] (708.0,70.0) -- (710.8,73.8);
  \draw[line width=5.0pt] (21.0,90.0) -- (23.3,57.7);
  \draw[line width=4.0pt] (23.3,57.7) -- (36.8,41.2);
  \draw[line width=5.0pt] (212.7,59.7) -- (213.0,90.0);
  \draw[line width=5.0pt] (57.0,119.0) -- (39.0,110.0);
  \draw[line width=4.0pt] (213.0,90.0) -- (195.0,108.0);
  \draw[line width=10.5pt] (604.0,248.0) -- (594.0,244.0);
  \draw[line width=5.0pt] (812.0,122.0) -- (817.0,116.0);
  \draw[line width=5.0pt] (22.0,91.0) -- (38.0,109.0);
  \draw[line width=5.0pt] (57.0,121.0) -- (56.0,187.0);
  \draw[line width=5.0pt] (163.0,192.0) -- (163.0,125.0);
  \draw[line width=7.7pt] (604.0,252.0) -- (602.0,259.0);
  \draw[line width=9.0pt] (693.0,69.0) -- (681.0,67.0);
  \draw[line width=4.0pt] (253.0,219.0) -- (306.0,254.0);
  \draw[line width=5.0pt] (441.0,222.0) -- (434.0,312.0);
  \draw[line width=5.0pt] (595.0,254.0) -- (603.0,251.0);
  \draw[line width=4.0pt] (434.0,315.0) -- (427.0,405.0);
  \draw[line width=4.0pt] (196.0,174.0) -- (195.0,110.0);
  \draw[line width=4.0pt] (354.0,352.0) -- (432.0,313.0);
  \draw[line width=4.0pt] (432.0,313.0) -- (361.0,263.0);
  \draw[line width=3.0pt] (569.0,269.0) -- (586.0,260.0);
  \draw[line width=7.7pt] (306.0,247.0) -- (306.0,254.0);
  \draw[line width=3.0pt] (544.0,282.0) -- (561.0,273.0);
  \draw[line width=4.0pt] (163.0,123.0) -- (194.0,109.0);
  \draw[line width=4.0pt] (516.0,273.0) -- (436.0,313.0);
  \draw[line width=6.5pt] (625.0,238.0) -- (605.0,248.0);
  \draw[line width=5.0pt] (107.0,131.0) -- (107.0,201.0);
  \draw[line width=7.0pt] (320.0,263.0) -- (307.0,256.0);
  \draw[line width=5.0pt] (1032.0,124.0) -- (991.6,72.6);
  \draw[line width=5.0pt] (674.2,21.0) -- (643.0,30.0);
  \draw[line width=4.0pt] (643.0,30.0) -- (608.5,50.5);
  \draw[line width=5.0pt] (797.5,268.0) -- (838.1,281.1);
  \draw[line width=4.0pt] (838.1,281.1) -- (859.0,291.0);

  %% ════ 曲线（prims2 的 curves —— 三段式，坐标同为 (y,x)）════
  \draw[line width=5.0pt] (775.8,220.8) .. controls (818.6,232.0) and (871.9,224.1) .. (897.8,184.2);
  \draw[line width=5.0pt] (898.0,174.1) .. controls (889.8,162.1) and (874.0,169.4) .. (862.3,166.0);
  \draw[line width=5.0pt] (862.3,166.0) .. controls (856.3,163.6) and (852.6,158.5) .. (854.0,151.8);
  \draw[line width=5.0pt] (854.0,151.8) .. controls (863.8,136.3) and (887.9,137.3) .. (894.0,118.7);
  \draw[line width=5.0pt] (894.0,118.7) .. controls (894.0,103.6) and (875.4,101.0) .. (864.0,99.0);
  \draw[line width=5.0pt] (164.0,193.0) .. controls (175.7,190.6) and (188.7,186.0) .. (196.0,176.0);
  \draw[line width=5.0pt] (711.4,94.4) .. controls (735.7,122.1) and (751.0,155.5) .. (758.0,190.8);
  \draw[line width=4.0pt] (197.0,175.0) .. controls (225.4,160.0) and (209.2,121.3) .. (215.0,95.2);
  \draw[line width=5.0pt] (710.2,68.0) .. controls (762.1,69.5) and (811.6,78.8) .. (859.2,94.2);
  \draw[line width=5.0pt] (710.8,73.8) .. controls (747.8,83.1) and (784.0,94.4) .. (817.0,115.0);
  \draw[line width=5.0pt] (36.8,41.2) .. controls (87.6,7.3) and (174.7,5.4) .. (212.7,59.7);
  \draw[line width=5.0pt] (56.0,188.0) .. controls (69.1,199.0) and (88.6,200.2) .. (106.0,202.0);
  \draw[line width=4.0pt] (788.0,168.0) .. controls (810.0,161.6) and (791.9,132.6) .. (812.0,122.0);
  \draw[line width=5.0pt] (163.0,193.0) .. controls (146.5,200.9) and (126.7,200.6) .. (108.0,202.0);
  \draw[line width=5.0pt] (180.0,104.0) .. controls (140.7,85.5) and (88.9,87.0) .. (50.0,105.0);
  \draw[line width=4.0pt] (35.0,171.0) .. controls (9.3,156.2) and (21.1,117.7) .. (20.0,92.0);
  \draw[line width=5.0pt] (55.0,188.0) .. controls (47.5,186.1) and (38.2,180.0) .. (36.0,172.0);
  \draw[line width=5.0pt] (787.0,168.0) .. controls (771.0,135.8) and (742.5,105.7) .. (711.0,88.0);
  \draw[line width=5.0pt] (1042.0,201.0) .. controls (1043.1,175.2) and (1039.9,148.0) .. (1032.0,124.0);
  \draw[line width=5.0pt] (991.6,72.6) .. controls (901.9,8.7) and (781.9,3.9) .. (674.2,21.0);
  \draw[line width=4.0pt] (608.5,50.5) .. controls (593.1,65.8) and (585.7,86.0) .. (581.0,106.2);
  \draw[line width=5.0pt] (581.0,106.2) .. controls (578.4,128.8) and (577.2,153.7) .. (585.0,174.9);
  \draw[line width=5.0pt] (585.0,174.9) .. controls (632.0,248.8) and (724.1,246.9) .. (797.5,268.0);
  \draw[line width=5.0pt] (859.0,291.0) .. controls (908.8,324.9) and (985.2,312.8) .. (1019.1,261.9);
  \draw[line width=5.0pt] (1019.1,261.9) .. controls (1038.1,247.5) and (1041.5,223.7) .. (1042.0,201.0);

  %% ════ 圆 / 椭圆（probe 精测）════
  \draw (435.2,313.4) circle (92.7);

  %% ════ 实心点 ════
  \fill (471.4,92.5) circle (3.7);
  \fill (301.9,257.9) circle (3.0);

  %% ════ 标签（probe 直接给的 \node 行）════
  %% ⚠ 全部补了 fill=white：文字在最上层，否则与曲线交叉干扰
     \node[anchor=center,inner sep=0pt,fill=white,font=\fontsize{30.7}{38.4}\selectfont] at (406.5,62.0) {\textsf{\textbf{\textit{R}}}};   % box(396,50,21,22)
     \node[anchor=center,inner sep=0pt,fill=white,font=\fontsize{22.2}{27.8}\selectfont] at (934.1,221.4) {\textsf{\textbf{1}}};   % box(928,213,7,16)
     \node[anchor=center,inner sep=0pt,fill=white,font=\fontsize{36.5}{45.6}\selectfont] at (919.4,229.9) {\textsf{\textbf{6}}};   % box(909,217,19,24)
     \node[anchor=center,inner sep=0pt,fill=white,font=\fontsize{30.7}{38.3}\selectfont] at (885.8,230.8) {\textsf{\textbf{\textit{b}}}};   % box(876,218,17,23)
     \node[anchor=center,inner sep=0pt,fill=white,font=\fontsize{30.7}{38.4}\selectfont] at (899.5,232.5) {\textsf{\textbf{(}}};   % box(895,218,8,28)
     \node[anchor=center,inner sep=0pt,fill=white,font=\fontsize{30.7}{38.4}\selectfont] at (946.5,233.5) {\textsf{\textbf{)}}};   % box(941,219,8,28)
     \node[anchor=center,inner sep=0pt,fill=white,font=\fontsize{22.2}{27.8}\selectfont] at (59.1,238.4) {\textsf{\textbf{1}}};   % box(53,230,7,16)
     \node[anchor=center,inner sep=0pt,fill=white,font=\fontsize{30.2}{37.8}\selectfont] at (44.3,247.0) {\textsf{\textbf{\textit{B}}}};   % box(33,235,19,23)
     \node[anchor=center,inner sep=0pt,fill=white,font=\fontsize{28.9}{36.1}\selectfont] at (106.9,247.5) {\textsf{\textbf{)}}};   % box(101,236,9,22)
     \node[anchor=center,inner sep=0pt,fill=white,font=\fontsize{36.0}{45.0}\selectfont] at (81.1,250.7) {\textsf{\textbf{×}}};   % box(72,240,17,16)
     \node[anchor=center,inner sep=0pt,fill=white,font=\fontsize{38.1}{47.6}\selectfont] at (1048.8,266.5) {\textsf{\textbf{\textit{x}}}};   % box(1037,253,21,22)
     \node[anchor=center,inner sep=0pt,fill=white,font=\fontsize{29.4}{36.8}\selectfont] at (261.0,268.5) {\textsf{\textbf{\textit{π}}}};   % box(253,258,17,18)
     \node[anchor=center,inner sep=0pt,fill=white,font=\fontsize{30.3}{37.9}\selectfont] at (602.8,291.5) {\textsf{\textbf{\textit{k}}}};   % box(594,279,17,22)
     \node[anchor=center,inner sep=0pt,fill=white,font=\fontsize{24.9}{31.2}\selectfont] at (527.1,407.4) {\textsf{\textbf{2}}};   % box(520,399,13,16)
     \node[anchor=center,inner sep=0pt,fill=white,font=\fontsize{28.9}{36.1}\selectfont] at (506.2,416.0) {\textsf{\textbf{\textit{B}}}};   % box(495,405,19,21)

  % 钉死画布：否则超出的图元/控制点会把页面撑大，1:1 回比就失准
  \pgfresetboundingbox
  \path[use as bounding box] (0,0) rectangle (1067,435);
\end{tikzpicture}
\end{document}

% ⚠ 待核清单（probe 拿不准，人眼过一遍）：
%   · 未识别/跳过：⑤ 跳过/未识别的字形
